\documentclass[11pt]{article}
\usepackage[margin=1in]{geometry}
\usepackage{times}
\usepackage{hyperref}
\hypersetup{colorlinks=true, linkcolor=black, urlcolor=blue, citecolor=black}
\title{Conservation Is Not Optional: The EmDrive and the Persistence of Reactionless Thrust Claims}
\author{Flyxion \\ \small Independent Researcher}
\date{}
\begin{document}
\maketitle

\begin{abstract}
The EmDrive, a truncated conical microwave resonator claimed to produce net thrust without expelling propellant, has been reported by several laboratories to generate small forces consistent with its predicted operation, and has correspondingly been proposed as evidence that momentum conservation admits exceptions under the right cavity geometry. This paper argues that momentum conservation is not a phenomenological regularity awaiting a documented counterexample but a consequence of translational symmetry with a depth of confirmation that a millinewton-scale laboratory anomaly cannot outweigh, and that the measured thrust is far more plausibly explained by mundane interactions between the test apparatus and its environment than by a mechanism with no theoretical channel to act through.
\end{abstract}

\section{Introduction}

Roger Shawyer's EmDrive proposes that microwaves bounced within an asymmetric, closed, tapered cavity can produce a net thrust in the direction of the wider end, without any propellant leaving the device. Several laboratories, including a NASA Eagleworks group, have reported measuring small forces, on the order of tens to hundreds of micronewtons, consistent in sign and rough magnitude with the device's predicted thrust direction. These results have been cited as evidence that the underlying theory, and by extension the conservation of momentum itself, requires revision.

\section{Why the Bar Here Is Unusually High}

Momentum conservation is not an empirical generalization discovered by repeated observation and subject to revision the way, say, a rate constant in chemistry might be. It follows, via Noether's theorem, from the translational symmetry of space, the fact that the laws of physics do not depend on where in space an experiment is performed. This symmetry has been tested to extraordinary precision across an enormous range of scales, from particle collider experiments to planetary orbital mechanics to the propagation of light across cosmological distances, and a violation of it would not merely require a new force or a correction term; it would require that space itself is not translationally uniform in a way that has escaped every prior experiment ever conducted at any other scale, but happens to manifest in a copper cavity roughly the size of a wastebasket. A micronewton-scale anomaly, measured at the very edge of a lab's noise floor, is not remotely sufficient evidence to overturn a symmetry this deeply and repeatedly confirmed, and treating it as sufficient inverts the appropriate burden of proof.

\section{What the Entropic Account of Gravity Adds}

Framing gravitation as the large-scale geometric bookkeeping of an ordered configuration rather than a force transmitted by a mediating field does not, by itself, forbid reactionless thrust, since the EmDrive's claimed mechanism is electromagnetic rather than gravitational. What it does is remove one of the more common post-hoc rescues offered for the device, the suggestion that it somehow pushes against the "quantum vacuum" or against spacetime itself in a manner analogous to how a propeller pushes against water. On the account assumed here, spacetime geometry is not a substrate with local momentum available to be borrowed from; it is the emergent bookkeeping of an entropy gradient at scales where the device in question is entirely negligible. There is, in short, nothing there for a microwave cavity to push against, which closes off the vacuum-drag rescue at the same time that Noether's theorem closes off the more direct claim.

\section{The Mundane Explanation}

Subsequent, more carefully controlled experiments, including a widely cited null result from a German group at TU Dresden using improved isolation from thermal expansion, magnetic interference from power cables, and interaction with Earth's magnetic field, found that the measured thrust tracked these mundane systematic effects rather than the cavity's resonant properties, and that a properly isolated cavity produced no anomalous force at all. Thermal expansion of the mounting apparatus during power-up, in particular, produces a force signature that mimics genuine thrust in exactly the timescale and magnitude range the earlier positive results reported.

\section{Objections}

A common reply holds that a null result from one lab does not retroactively invalidate positive results from several others, and that a systematic replication effort across differently designed test stands would settle the question more conclusively than any single refutation. This is true as a general methodological point, but it also describes exactly what has occurred: as thrust-balance isolation techniques have improved across the field, the measured effect has shrunk toward zero rather than converging on a stable, well-characterized nonzero value, which is the signature of instrumental artifact being progressively eliminated rather than of a real effect being progressively resolved.

\section{Conclusion}

A microwave cavity with a millinewton-scale reported anomaly is not remotely sufficient grounds to revise a symmetry as deeply confirmed as momentum conservation, and the entropic account of gravitation removes even the fallback appeal to a vacuum substrate the device might be pushing against. The improving trend toward null results as measurement technique has matured is exactly what the mundane explanation predicts and exactly what the reactionless-thrust explanation does not.

\end{document}
